\begin{scanpage}{0088}{71}
which corresponds to the physical optics approach [Equation (7.20)].

The functions $\overline\Sigma$ and $\Sigma$ allow one to calculate
\[\sigma_E=\pi a^2|\overline\Sigma|^2,\quad\sigma_H=\pi a^2|\Sigma|^2\tag{12.17}\]
which are the effective scattering surfaces with the E- and H-polarizations of the incident wave. Let us recall that, by definition, the effective scattering surface is a quantity equal to
\[\sigma=4\pi R^2\frac{|\mathbf S|}{|\mathbf S_0|},\tag{12.18}\]
where
\[\mathbf S=\frac{c}{8\pi}\operatorname{Re}[\mathbf E\mathbf H^*]\tag{12.19}\]
which is the energy flux density averaged over one oscillation cycle (the Poynting vector) in the scattered wave, and $\mathbf S_0$ is a similar quantity for the incident wave.

In this way, we obtained the expressions for the fringing field which approximately take into account the nonuniform part of the current. In the incident plane $(\varphi=\pm\frac\pi2)$, they have a form which is rather simple and convenient for calculations. It is also interesting that in this case they satisfy the reciprocity principle as distinct from expressions (10.03) and (10.04) which correspond to the uniform part of the current. It is not difficult to prove this by verifying that Equations (12.11) are not changed with the simultaneous replacement of $\gamma$ by $\vartheta$ and of $\vartheta$ by $\gamma$, and Equations (12.12) are not changed with the replacement of $\vartheta$ by $\pi-\gamma$ and of $\gamma$ by $\pi-\vartheta$ [in the case $\varphi=-\frac\pi2$].

However, Equations (12.11) and (12.12) lead to a discontinuity of the magnetic field tangential component $H_\varphi$ on the plane $z=0$ in which the disk lies. As in the case of diffraction by a strip, the
\end{scanpage}
